% lecture03.tex - 线性变换与矩阵
% 第 3 章: 线性变换与矩阵

\chapter{线性变换与矩阵}
\label{ch:transformations}

\begin{abstract}
线性变换 (linear transformation) 是线性代数的核心。在本章中, 我们将学习任何线性变换都可以用一个矩阵 (matrix) 来表示, 以及矩阵乘法与变换复合之间的关系。
\end{abstract}

% ============================================
\section{什么是变换?}
% ============================================

\begin{definition}[变换 (transformation)]
\vocab{变换}是一个把向量变成另一个向量的函数:
\[
T: \Rtwo \to \Rtwo
\]
每个输入向量都被映射到一个输出向量。
\end{definition}

\begin{intuition}
我们用``变换''而不用``函数'', 是因为我们想强调\textbf{运动}。想象空间中的每个向量都被``移动''到了新的位置。

再想象坐标系的网格线。变换会拉伸、压缩、旋转这些网格, 或以其他方式改变它们。
\end{intuition}

% ============================================
\section{线性变换}
% ============================================

\begin{definition}[线性变换 (linear transformation)]
变换 $T$ 是\vocab{线性变换}, 当且仅当以下两个条件成立:
\begin{enumerate}
    \item \textbf{可加性:} $T(\vv + \vw) = T(\vv) + T(\vw)$
    \item \textbf{齐次性:} $T(c\vv) = c \cdot T(\vv)$
\end{enumerate}
其中 $\vv, \vw$ 为任意向量, $c$ 为任意标量。
\end{definition}

\begin{intuition}
线性变换可以从两个几何性质来识别:
\begin{enumerate}
    \item \textbf{直线变换后仍是直线} (不会变弯)
    \item \textbf{原点保持不动}
\end{enumerate}

如果变换后网格线仍然平行且等距, 那它就是线性的。
\end{intuition}

\begin{center}
\begin{tikzpicture}[scale=0.8]
    % 原网格
    \begin{scope}[shift={(0,0)}]
        \draw[grid] (-2,-2) grid (2,2);
        \draw[axis] (-2,0) -- (2,0);
        \draw[axis] (0,-2) -- (0,2);
        \draw[basis, -{Stealth[length=3mm]}] (0,0) -- (1,0) node[below] {$\vi$};
        \draw[basis, -{Stealth[length=3mm]}] (0,0) -- (0,1) node[left] {$\vj$};
        \node at (0,-3) {变换前};
    \end{scope}

    % 箭头
    \draw[-{Stealth[length=5mm]}, very thick] (3,0) -- (5,0) node[midway, above] {$T$};

    % 变换后的网格 (错切)
    \begin{scope}[shift={(8,0)}]
        \draw[blue!30, thin] (-2,-2) -- (0,2) -- (4,2) -- (2,-2) -- cycle;
        \foreach \y in {-1,0,1} {
            \draw[blue!30, thin] ({-2+\y},-2) -- ({\y+2},2);
        }
        \foreach \x in {-2,-1,0,1,2} {
            \draw[blue!30, thin] ({\x-1},-2) -- ({\x+1},2);
        }
        \draw[axis] (-2,0) -- (4,0);
        \draw[axis] (0,-2) -- (2,2);
        \draw[transformed, -{Stealth[length=3mm]}, very thick] (0,0) -- (1,0) node[below] {$T(\vi)$};
        \draw[transformed, -{Stealth[length=3mm]}, very thick] (0,0) -- (1,1) node[above left] {$T(\vj)$};
        \node at (1,-3) {变换后};
    \end{scope}
\end{tikzpicture}
\end{center}

% ============================================
\section{线性变换的矩阵}
% ============================================

\begin{theorem}[基本定理]
每个线性变换 $T: \Rtwo \to \Rtwo$ 完全由 $T$ 把基向量 (basis) $\vi$ 与 $\vj$ 送到哪里决定。
\end{theorem}

\begin{proof}
任何向量 $\vv = \twovec{x}{y}$ 都可以写成: $\vv = x\vi + y\vj$

利用线性性质:
\[
T(\vv) = T(x\vi + y\vj) = xT(\vi) + yT(\vj)
\]
因此只要知道 $T(\vi)$ 与 $T(\vj)$, 就能算出任何 $\vv$ 的 $T(\vv)$。
\end{proof}

\begin{definition}[变换矩阵 (transformation matrix)]
若 $T(\vi) = \twovec{a}{c}$, $T(\vj) = \twovec{b}{d}$, 则 $T$ 的\vocab{变换矩阵}为:
\[
\mA = \twomat{a}{b}{c}{d}
\]
第一列是 $\vi$ 的落点, 第二列是 $\vj$ 的落点。
\end{definition}

\begin{important}
\textbf{黄金法则:} 矩阵的列向量 (column vector) = 基向量的落点
\end{important}

% ============================================
\section{矩阵与向量相乘}
% ============================================

\begin{definition}[矩阵乘以向量]
\[
\twomat{a}{b}{c}{d} \twovec{x}{y} = x\twovec{a}{c} + y\twovec{b}{d} = \twovec{ax + by}{cx + dy}
\]
\end{definition}

\begin{intuition}
矩阵乘以向量意味着:
\begin{enumerate}
    \item 把输入的 $x$ 分量乘以第一列
    \item 把输入的 $y$ 分量乘以第二列
    \item 把结果加起来
\end{enumerate}
这正是新基向量的线性组合!
\end{intuition}

\begin{example}
逆时针旋转 $90°$:
\[
T(\vi) = \twovec{0}{1}, \quad T(\vj) = \twovec{-1}{0}
\]
于是旋转矩阵为:
\[
\mR = \twomat{0}{-1}{1}{0}
\]
验证: $\mR\twovec{1}{0} = \twovec{0}{1}$ \checkmark
\end{example}

% ============================================
\section{重要变换举例}
% ============================================

\subsection{旋转}

\begin{definition}[旋转矩阵]
逆时针旋转 $\theta$ 角:
\[
\mR_\theta = \twomat{\cos\theta}{-\sin\theta}{\sin\theta}{\cos\theta}
\]
\end{definition}

\begin{example}
旋转 $45°$:
\[
\mR_{45°} = \twomat{\frac{\sqrt{2}}{2}}{-\frac{\sqrt{2}}{2}}{\frac{\sqrt{2}}{2}}{\frac{\sqrt{2}}{2}}
\]
\end{example}

\subsection{缩放}

\begin{definition}[缩放矩阵]
沿 $x$ 轴缩放 $s_x$ 倍、沿 $y$ 轴缩放 $s_y$ 倍:
\[
\mS = \twomat{s_x}{0}{0}{s_y}
\]
\end{definition}

\begin{practical}
\textbf{在计算机图形学中的应用:} 要放大或缩小一幅图像:
\begin{itemize}
    \item $s_x = s_y = 2$: 图像变为原来的两倍
    \item $s_x = 1, s_y = 0.5$: 图像在竖直方向被压缩
\end{itemize}
\end{practical}

\subsection{错切}

\begin{definition}[水平错切矩阵]
\[
\mH = \twomat{1}{k}{0}{1}
\]
其中 $k$ 为错切量。
\end{definition}

\begin{center}
\begin{tikzpicture}[scale=1]
    % 原图形
    \begin{scope}[shift={(0,0)}]
        \draw[fill=blue!20] (0,0) -- (1,0) -- (1,1) -- (0,1) -- cycle;
        \draw[basis] (0,0) -- (1,0) node[midway, below] {$\vi$};
        \draw[basis] (0,0) -- (0,1) node[midway, left] {$\vj$};
        \node at (0.5,-0.7) {原图形};
    \end{scope}

    % 错切后
    \begin{scope}[shift={(4,0)}]
        \draw[fill=blue!20] (0,0) -- (1,0) -- (2,1) -- (1,1) -- cycle;
        \draw[transformed, very thick] (0,0) -- (1,0) node[midway, below] {$T(\vi)$};
        \draw[transformed, very thick] (0,0) -- (1,1) node[midway, left] {$T(\vj)$};
        \node at (1,-0.7) {错切 ($k=1$)};
    \end{scope}
\end{tikzpicture}
\end{center}

\subsection{反射}

\begin{example}
关于 $x$ 轴的反射:
\[
\twomat{1}{0}{0}{-1}
\]

关于直线 $y = x$ 的反射:
\[
\twomat{0}{1}{1}{0}
\]
\end{example}

% ============================================
\section{矩阵乘法即变换的复合}
% ============================================

\begin{theorem}[变换的复合]
若 $\mA$ 是变换 $T_1$ 的矩阵, $\mB$ 是变换 $T_2$ 的矩阵, 则:
\[
\mB \cdot \mA = \text{矩阵 } (T_2 \circ T_1)
\]
即先做 $T_1$, 再做 $T_2$。
\end{theorem}

\begin{warning}
顺序很重要! 一般来说 $\mA\mB \neq \mB\mA$。

先旋转再错切 $\neq$ 先错切再旋转
\end{warning}

\begin{definition}[矩阵乘法 (matrix multiplication)]
\[
\twomat{a}{b}{c}{d} \twomat{e}{f}{g}{h} = \twomat{ae+bg}{af+bh}{ce+dg}{cf+dh}
\]
\end{definition}

\begin{intuition}
计算 $\mB\mA$ 的列:
\begin{itemize}
    \item 第一列: $\mB$ 乘以 $\mA$ 的第一列 = $\vi$ 经过两次变换后的落点
    \item 第二列: $\mB$ 乘以 $\mA$ 的第二列 = $\vj$ 经过两次变换后的落点
\end{itemize}
\end{intuition}

\begin{example}
旋转 $90°$ 后再错切:
\[
\underbrace{\twomat{1}{1}{0}{1}}_{\text{错切}} \underbrace{\twomat{0}{-1}{1}{0}}_{\text{旋转}} = \twomat{1}{-1}{1}{0}
\]

验证: $\vi \to \twovec{0}{1} \to \twovec{1}{1}$ \checkmark
\end{example}

% ============================================
\section{实际应用}
% ============================================

\begin{practical}
\textbf{计算机图形学与电子游戏}

游戏中的每个物体都由一组向量 (顶点) 描述。要移动、旋转或缩放一个物体, 只需把所有顶点乘以相应的矩阵。

一段动画 = 一连串矩阵乘法
\end{practical}

\begin{practical}
\textbf{图像处理}

模糊、锐化、边缘检测等图像滤波器都是用矩阵乘法实现的。
\end{practical}

% ============================================
\section{练习}
% ============================================

\begin{exercise}
写出把 $\vi$ 送到 $\twovec{2}{1}$、把 $\vj$ 送到 $\twovec{-1}{3}$ 的变换的矩阵。
\end{exercise}

\begin{exercise}
计算以下乘积:
\[
\twomat{1}{2}{3}{4} \twovec{5}{6}
\]
\end{exercise}

\begin{exercise}
写出旋转 $180°$ 的矩阵, 并证明它等于 $-\mI$。
\end{exercise}

\begin{exercise}
把以下两个矩阵按两种顺序相乘, 并证明 $\mA\mB \neq \mB\mA$:
\[
\mA = \twomat{1}{2}{0}{1}, \quad \mB = \twomat{0}{1}{1}{0}
\]
\end{exercise}

\begin{exercise}[挑战]
求关于直线 $y = 2x$ 的反射矩阵。

\textit{提示: 这条直线与 $x$ 轴的夹角为 $\arctan(2)$。}
\end{exercise}

\begin{problem}
证明: 旋转 $\alpha$ 角后再旋转 $\beta$ 角, 等于旋转 $\alpha + \beta$ 角。
\end{problem}

\begin{problem}
若 $\mA^2 = \mA$ (即 $\mA$ 为投影矩阵), 它的几何意义是什么?
\end{problem}
